\chapter{Анализ результатов}
\label{ch:6}

В данной главе представлен сводный анализ экспериментальных результатов разработанных алгоритмов укладки.

\section{Сравнение алгоритмов для одной паллеты}

Тестирование проводилось на наборе из 436 реальных задач. Размеры паллеты: $1200 \times 800$ мм, нормировочная высота: 2200 мм. Временной лимит 5~с для метаэвристических алгоритмов, веса целевой функции $w_1 = 1{,}0$, $w_2 = 2{,}0$. Сводные численные показатели приведены в таблице~\ref{tab:evyukh_all_solvers} (раздел~\ref{sec:evyukh_experiments}).

Алгоритмы разделяются на две группы по критерию наличия ограничения на устойчивость. Без ограничения (DBLFSolver, FFDSolver) достигается наивысшая перколяция (65--66\%), но минимальный коэффициент опоры падает ниже 3\%. Метаэвристические алгоритмы (LNSSolver, GeneticSolver) обеспечивают баланс: перколяция $\approx$63,5\% при минимальном коэффициенте опоры свыше 62\%.

LNSSolver является лучшим алгоритмом по совокупности метрик при ограничении на время порядка 5~с: он превосходит GreedySolver на 7,45~п.п. по перколяции и на 34,43~п.п. по минимальному коэффициенту опоры, а также немного опережает GeneticSolver по обоим показателям.

\section{Результаты в мультипаллетном режиме}

Мультипаллетный режим тестировался на том же наборе тестов, объединённых в группы по нескольку заказов. Сводные показатели --- в таблице~\ref{tab:evyukh_multi} (раздел~\ref{sec:evyukh_experiments}).

В мультипаллетном режиме FFDSolver показывает критически низкий баланс высот (52\%), поскольку крупные коробки концентрируются на первых паллетах. LNSSolver демонстрирует лучшую перколяцию среди метаэвристик (62,09\%), сохраняя высокий баланс высот (84,26\%) и приемлемый коэффициент опоры.